%****************************************************
\subsection{Behavioral vs. Structural Descriptions}
%***************************************************
%***************************************
\subsubsection{Circuit Characterization}
%***************************************

In a digital system we can distinguish two extreme types of circuits:
\begin{itemize}
\item purely logical circuits that operate on some binary signals with some symbolic meaning (example: the transcoder for displaying the numbers from 0 to 9 using 7 segments, represented in Figure \ref{fig2}a)

\item purely numerical circuits that operate on binary signals with numerical significance (example: a modulo 16 adder for two numbers represented by 4 bits each, represented in Figure \ref{fig2}b)
\end{itemize}
Obviously, an intermediate version is always possible (example: a circuit that compares two numbers, A and B, and indicates through three logical signals of one bit each: $A>B, \; A=B, \; A\leq B$)

\placedrawing[H]{figures/01_03_7seg_and_adder.lp}{The version of digital circuits. {\bf a.} Logic circuit: trans-coder for seven-segment display. {\bf b.} Numeric circuit: four-bit numbers adder.}{fig2}


The description of a digital circuit can be done in two ways:
\begin{itemize}
\item  by describing the functional behavior
\item  by describing the internal structure of the circuit
\end{itemize}
We exemplify the two methods in the case of the addition function.


\begin{exemplu}\label{adder}
Let us use the {\em System Verilog} HDL to describe an adder for 4-bit
numbers (see Figure \ref{first_ex}a). The description which
follows is a {\bf behavioral} one, because we know {\bf what} we
intend to design, but we do not know yet {\bf how} to design the
internal structure of an adder.

\placedrawing[h]{figures/01_04_4bit_adder.lp}{{\bf The first examples of  digital
systems.} {\small {\bf a.} The two 4-bit numbers adder, called
{\tt adder}. {\bf b.} The structure of an adder for 3 4-bit
numbers, called {\tt threeAdder}.} }{first_ex}

The Verilog code describing the module {\tt adder} is:
{\small
{\em
\begin{shaded*}
\begin{lstlisting}[language=Verilog, caption= , morekeywords={always_ff, always_comb, logic, loop, repeat, doInParallel, for}]
/*************************************************************************
File name:      adder.sv
Circuit name:   Adder
Description:    The module 'adder' has 2 4-bit inputs and one 4-bit output
                The circuit adds modulo 16; do not use or provide carry
                signal
*************************************************************************/
module adder(   output logic [3:0] out ,    // 4-bit output
                input  logic [3:0] in0 ,    // 4-bit input
                input  logic [3:0] in1 );   // 4-bit input

    assign out = in0 + in1;
endmodule
\end{lstlisting}
\end{shaded*}
}
}

The story just told by the previous Verilog module is: ``the 4-bit adder has two inputs, {\tt in0, in1}, one output, {\tt out}, and
its output is continuously assigned to the value obtained by
adding modulo 16 the two input numbers".

  $\diamond$
\end{exemplu}

What we just learned from the previous first simple example is
summarized in the following {\bf SystemVerilogSummary}.

\begin{verisum}{\em :

\hspace{-0.8cm} \rule[5mm]{16cm}{0.25mm} \vspace{-0.7cm}

\begin{description}
\item[/*]: begin comment
\item[*/]: end comment
\item[module]: keyword which indicates the beginning of the
description of a circuit as a module having the name which
immediately follows (in our example, the name is: {\tt adder})

\item[endmodule]: keyword which indicates the end of the module's
description which started with the previous keyword {\bf module}

\item[output]: keyword used to declare a terminal as an output (in
our example the terminal {\tt out} is declared as output)

\item[input]: keyword used to declare the terminal as an input (in
our example the terminals {\tt in0} and {\tt in1} are declared as
inputs)

\item[logic]: a data type with 4-state bits: {\tt 0, 1, x, z} where 0 means low, 1 means high, x means unknown, and z means an undriven net.

\item[initial]: a block which starts at the beginning of simulation

\item[begin ... end]: block delimiters

\item[\#]: indicate a raw value

\item[forever]: indicate an unending operation

\item[$\sim$]: negation operator

\item[assign]: keyword called the {\em continuous assignment},
used here to specify the function performed by the module (the
output {\tt out} takes continuously the value computed by adding
the two input numbers)

\item[(...)]: delimiters used to delimit the list of terminals
(external connections)

\item[,]: delimiter to separate each terminal within a list of
terminals

\item[;]: delimiter for end of line

\item[[ $ \dots$ ]]: delimiters which contains the definition of
the bits associated with a connection, for example [3:0] define
the number of bits for the three connections in the previous
example

\item[+]: the operator {\tt add}, the only one used in the
previous example.
\end{description}
\rule[5mm]{16cm}{0.25mm}
}
\end{verisum}

The description of a digital system is a {\bf hierarchical
construct} starting from a top module populated by modules, which
are similarly defined. The process continues until very simple
module are directly described. Thus, the functions $f$ and $g$ are specified by HDL programs
(in our case, in the previous example by a Verilog program).

The main characteristic of the digital design is {\bf modularity}. A
problem is decomposed in many simpler problems, which are solved
similarly, and so on until very simple problems are identified.
Modularity means also to define as many as possible identical
modules in each design. This allow to replicate many times the
same module, already designed and validated. {\em Many \& simple}
modules! Is the main slogan of the digital designer. Let's take
another example which uses as module the one just defined in the
previous example.

\begin{exemplu}
The previously exemplified module ({\tt adder}) will be used to
design a modulo 16 3-number adder, called {\tt threeAdder} (see Figure
\ref{first_ex}b). It adds 3 4-bit numbers providing a 4-bit result
(modulo 16 sum). Follows the {\bf structural} description:
{\small
{\em
\begin{shaded*}
\begin{lstlisting}[language=Verilog, caption= , morekeywords={always_ff, always_comb, logic, loop, repeat, doInParallel, for}]
/*************************************************************************
File name:      threeAdder.sv
Circuit name:   Three Input Adder
Description:    The module 'threeAdder' has 3 4-bit inputs and one 4-bit
                output. The circuit adds modulo 16 three numbers;
                do not provide carry output
*************************************************************************/
module threeAdder( output   logic [3:0]  out,
                   input    logic [3:0]  in0,
                   input    logic [3:0]  in1,
                   input    logic [3:0]  in2);

    logic [3:0] sum ;

    adder   inAdder(.out(sum),
                    .in0(in1),
                    .in1(in2)),
            outAdder(.out(out),
                     .in0(in0),
                     .in1(sum));
endmodule
\end{lstlisting}
\end{shaded*}
}
}

Two modules of {\tt adder} type (defined in the previous example)
are instantiated as {\tt inAdder}, {\tt outAdder}, they are
interconnected using the wire {\tt sum}, and are connected to the
terminals of the {\tt threeAdder} module. The resulting structure
computes the sum of three numbers. $\diamond$
\end{exemplu}

\begin{verisum}{\em :

\hspace{-0.8cm} \rule[5mm]{16cm}{0.25mm} \vspace{-0.7cm}

\begin{itemize}
\item How a previously defined module (in our example: {\tt
adder}) is two times instantiated using two different names ({\tt
inAdder} and {\tt outAdder} in our example)

\item A ``safe" way to allocate the terminals for a module
previously defined and instantiated inside the current module:
each original terminal name is preceded by a dot, and followed by
a parenthesis containing the name of the wire or of the terminal
where it is connected (in our example, {\tt outAdder( ...
.in1(sum))} means: the terminal {\tt in1} of the instance {\tt
outAdder} is connected to the wire {\tt sum})

\item The successive instantiations of the same module can be
separated by a ``,".
\end{itemize}
\rule[5mm]{16cm}{0.25mm} }
\end{verisum}

While the module {\tt adder} is a {\em behavioral} description,
the module {\tt threeAdder} is a {\em structural} one. The first
tells us {\em what} is the function of the module, and the second
tells us {\em how} its functionality is performed by using a
structure containing two instantiation of a previously defined
subsystems, and an internal connection.


%%%%%%%%%%%%%%%%%%%%%%%%%%%%%%%%%%%%%%%%%%%%%%%%%%%%%%%%%%%%%%%%%%%%
%******************************
\subsubsection{CMOS switches}
%******************************
%%%%%%%%%%%%%%%%%%%%%%%%%%%%%%%%%%%%%%%%%%%%%%%%%%%%%%%%%%%%%%%%%%%%%%%
A logic gates consists in a network of interconnected switches
implemented using the two type of MOS transistors: p-MOS and
n-MOS. How behaves the two type of transistors in specific
configurations is presented in Figure \ref{0_switch}.

A switch connected to $V_{DD}$ is implemented using a p-MOS
transistor. It is represented in Figure \ref{0_switch}a {\em off}
(generating $z$, which means {\em Hi-Z}: no signal, neither 0, nor
1) and in Figure \ref{0_switch}b it is represented {\em on}
(generating 1 logic, or {\em truth}).

A switch connected to {\em ground} is implemented using a n-MOS
transistor. It is represented in Figure \ref{0_switch}c {\em off}
(generating $z$, which means {\em Hi-Z}: no signal, neither 0, nor
1) and in Figure \ref{0_switch}e it is represented {\em on}
(generating 0 logic, or {\em false}).

\placedrawing[h]{figures/02_02_transistors.lp}{{\bf Basic switches.} {\small {\bf a}.
Open switch connected to $V_{DD}$. {\bf b}. Closed switch connected to
$V_{DD}$. {\bf c}. Open switch connected to {\em ground}. {\bf d}.
Closed switch connected to {\em ground}.}}{0_switch}

A MOS transistor works very well as an {\em on-off} switch
connecting its drain to a certain potential.  A p-MOS transistor
can be used to connect its drain to a high potential when its
gates is connected to ground, and an n-MOS transistor can connect
its drain to ground if its gates is connected to a high potential.
This complementary behavior is used to build the elementary logic
circuits.

In Figure \ref{src} is presented the {\em switch-resistor-capacitor
model} (SRC). If $V_{GS} < V_{T}$ then the transistor is {\bf off}, if
$V_{GS} \geq V_{T}$ then the transistor is {\bf on}. In both cases the
input of the transistor behaves like a capacitor, the gate-source
capacitor $C_{GS}$.

When the transistor is on its drain-source resistance is:
$$R_{ON} = R_{n}\frac{L}{W}$$
where: $L$ is the channel length, $W$ is the channel width, and $R_{n}$
is the resistance per square. The length $L$ is a constant
characterizing a certain technology. For example, if $L = 0.13\mu m$
this means it is about a $0.13\mu m$ process.

The input capacitor has the value:
$$C_{GS} = \frac{\epsilon_{OX} LW}{d}.$$
The value:
$$C_{OX} = \frac{\epsilon_{OX}}{d}$$
where: $\epsilon_{OX} \approx 3.9 \epsilon_{0}$ is the permittivity of
the silicon dioxide, is  the gate-to-channel capacitance per unit area
of the MOSFET gate.

In this conditions the gate input current is:
$$i_{G} = C_{GS} \frac{dv_{GS}}{dt}$$

\placedrawing[h]{figures/02_03_mosfet.lp}{{\bf The MOSFET switch.} {\small The {\em
switch-resistor-capacitor} model consists in the two states: OF
($V_{GS} < V_{T}$), and ON ($V_{GS} \geq V_{T}$). In both states the
input is defined by the capacitor $C_{GS}$.}}{src}


\begin{exemplu}
For an AND gate with low strength, with $W = 1.8\mu m$, in $0.13
\mu m$ technology, supposing $C_{OX} = 4fF/\mu m^{2}$, results the
input capacitance: $$C_{GS} = 4 \times 0.13 \times 1.8 fF = 0.936
fF$$ Assuming $R_{n} = 5K\Omega$, results for the same gate:
$$R_{ON} = 5 \times \frac{0.13}{1.8} \; K\Omega = 361 \Omega$$

$\diamond$
\end{exemplu}

\subsubsection{The Inverter}


\paragraph{The static behavior.}
The smallest and simplest logic circuit -- the invertor -- can be built
using a pair of complementary transistors, connecting together the two
gates as input and the two drains as output, while the n-MOS source is
connected to ground (interpreted as logic 0) and the p-MOS source to
$V_{DD}$ (interpreted as logic 1). Results the circuit represented in
Figure \ref{0_inv}.

\placedrawing[h]{figures/02_04_inverter.lp}{{\bf Building an invertor.} {\small {\bf
a}. The invertor circuit. {\bf b}. The logic symbol for the invertor
circuit.}}{0_inv}

The behavior of the invertor consist in combining the behaviors of the
two switches previously defined. For $in=0$ pMOS is $on$ and nMOS is
$off$ the output generating $V_{DD}$ which means 1. For $in=1$ pMOS is
$off$ and nMOS is $on$ the output generating 0.

The static behavior of the inverter (or NOT) circuit can be easy
explained starting from the switches described in Figure
\ref{0_switch}. Connecting together a switch generating $z$ with a
switch generating 1 or 0, the connection point will generate 0 or
1.

\paragraph{Dynamic behavior.}
The propagation time of an inverter can be analyzed using the two
serially connected invertors represented in Figure \ref{switch}. The
delay of the first invertor is generated by its capacitive load,
$C_{L}$, composed by:
\begin{itemize}
\item its parasitic drain/bulk capacitance, $C_{DB}$, is the intrinsic  output capacitance of the first invertor
\item wiring capacitance, $C_{wire}$, which depends on the length of the wire (of width $W_w$ and of length $L_w$) connected between the two invertors: $$C_{wire} = C_{thickox} W_w L_w$$
\item next stage input capacitance, $C_{G}$, approximated by summing the gate capacitance for pMOC and nMOS transistors: $$C_G = C_{Gp} + C_{Gn} = C_{ox}(W_p L_p + W_n L_n)$$
\end{itemize}

\placedrawing[h]{figures/02_05_propagation_time.lp}{{\bf The propagation time.} {\small
}}{switch}

The total load capacitance
$$C_L = C_{DB} + C_{wire} + C_G$$
is sometimes dominated by $C_{wire}$. For short connections $C_G$ dominates, while for big {\em fan-out}\index{fan-out: the number on input gates  connected to the output of a digital circuit} both, $C_{wire}$ and $C_G$ must be considered.

The signal $V_A$ is used to measure the propagation time of the first NOT in Figure \ref{switch}a. It is generated by an ideal pulse generator with output impedance 0. Thus, the rising time and the falling time of this signal are considered 0 (the input capacitance of the NOT circuit is charged  or discharged in no time).

The two delay times (see Figure \ref{switch}c) associated to an invertor
(to a gate in the general case) are defined as follows:
\begin{itemize}
    \item $t_{pLH}$: the time interval between the moment the input
    switches in 0 and the output reaches $V_{OH}/2$ coming from 0
    \item $t_{pHL}$: the time interval between the moment the input
    switches in 1 and the output reaches $V_{OH}/2$ coming from $V_{OH}$
\end{itemize}

Let us consider the transition of $V_A$ from 0 to $V_{OH}$ at $t_r$ (rise edge). Before transition, at $t_r^-$, $C_L$ is fully charged and $V_B = V_{OH}$. In Figure \ref{switch}b is represented the equivalent circuit at $t_r^+$, when pMOS is off and nMOS is on. In this moment starts the process of discharging  the capacitance $C_L$ at the constant current
$$I_{Dn(sat)} = \frac{1}{2} \mu_n C_{ox} \frac{W_n}{L_n} (V_{OH} - V_{Tn})^2$$
In Figure \ref{switch1}, at $t_r^-$ the transistor is cut, $I_{Dn} = 0$. At $t_r^+$ the  nMOS transistor switch in saturation and becomes an ideal {\em constant current generator} which starts to discharge $C_L$ linearly at the constant current $I_{Dn(sat)}$. The process continue until $V_{OUT} = V_{OH}$, according to the definition of $t_{pHL}$.


\placedrawing[h]{figures/02_06_nmos_output.lp}{{\bf The output characteristic of the nMOS transistor.} {\small
}}{switch1}

In order to compute $t_{pHL}$ we take into consideration the constant value of the discharging current which provide a linear variation of $v_{OUT}$.
$$\frac{dv_{out}}{dt} = \frac{d}{dt}(\frac{q_L}{C_L}) = \frac{-I_{Dn(sat)}}{C_L}$$
$$\frac{dv_{out}}{dt} = \frac{\frac{V_{OH}}{2} - V_{OH}}{t_{pHL}}$$
We solve the equations for $t_{pHL}$:
$$t_{pHL} = C_L \frac{1}{\mu_n C_{ox} \frac{W_n}{L_n} (V_{OH} - V_{Tn})} \frac{V_{OH}}{V_{OH} - V_{Tn}}$$
Because:
$$R_{ONn} = \frac{1}{\mu_n C_{ox} \frac{W_n}{L_n} (V_{OH} - V_{Tn})}$$
results:
$$t_{pHL} = C_L R_{ONn} \frac{1}{1 - \frac{V_{Tn}}{V_{OH}}} = k_n R_{ONn} C_L = k_n \tau_{nL} $$
where:
\begin{itemize}
\item $\tau_{nL}$ is the {\em constant time} associated to the H-L transition
\item $k_n$ is a constant associated to the technology we use; it goes down when $V_{OH}$ increases or $V_T$ decreases
\end{itemize}

The speed of a gate depends by its dimension and by the capacitive load
it drives. For a big $W$ the value of $R_{ON}$ is small charging or
discharging $C_{L}$ faster.

For $t_{pLH}$ the approach is similar. Results: $t_{pLH} = k_p \tau_{pL}$.

By definition the propagation time associated to a circuit is: $$t_p = (t_{pLH} + t_{pHL})/2$$
its value being dominated by the value of $C_L$ and the size (width) of the two transistors, $W_n$ and $W_p$.

